% Français 9117, batch 4 — Gallica views 61–80.
% Pass: Opus 5 (claude-opus-5), 2026-09-14 — first pass, unchecked against
% the views by a human.
%
% What the twenty views hold. One continuous stretch of a single prose work,
% the « Considérations générales » of the catalogue title: a working fair copy
% in Germain's hand, leaves written on both sides, lightly but constantly
% corrected (words struck in the line and rewritten above it, additions in
% the margin). No mathematics is set as formulae; the argument is about
% analogy, and mechanics is used as its instrument. Nineteen of the twenty
% views carry text; view 66 is the blank verso of the leaf written at view 65
% (only show-through and a struck number).
%
% View 61 opens mid-sentence (« qui a séduit autrefois l'esprit humain »),
% continuing view 60 of batch 3. View 80 breaks off mid-sentence (« où
% l'objet de »), continued in batch 5.
%
%   v. 61–63   The idea of being as the most abstract of our ideas; the
%              sentiment of analogy, born of the uniformity of the conditions
%              of being, as the source of both errors and happy thoughts; the
%              « esprit de système » (taking one particular truth as the base
%              of a whole order of facts, seeking identities where there are
%              only analogies) giving way to methodical research: experiment,
%              classification, laws, then theory — and the same integrals
%              representing unrelated orders of phenomena.
%   v. 64–65   The laws of being govern the intellectual order too; were the
%              language of calculation applicable to morals, politics,
%              metaphysics and taste, the likeness of the formulas would show
%              it. Example: perturbing causes die out quickly — in morals
%              (fraud), politics, reasoning (error), taste (fashion).
%   v. 67–72   Rational mechanics and political science compared at length:
%              equilibrium of opposed forces; stable and unstable equilibrium
%              and states that settle or are overturned; centre of gravity
%              and the tendency of the age (« l'opinion »); impulse through
%              the centre of gravity or off it (rotation, frontier provinces
%              as the most agitated points); hard, elastic and soft bodies
%              compared with men of interest, of passion, and of inertia.
%   v. 72–76   None of the three characters is ever pure; in troubled times
%              everyone becomes passionate, directions are uncertain, new men
%              of great natural force appear; « l'égalité est une erreur »,
%              illustrated by forces vives and the lever; perturbations are
%              functions of time.
%   v. 76–79   Analogies between the physical, moral and intellectual orders
%              as attested by figurative language (« monstre », and what
%              anatomy later taught about monstrosity); one well-established
%              relation between two subjects announces many others; analogy
%              is not a « méthode » of invention but exists by itself; the
%              tendency to generalise precipitated judgement ahead of
%              experience.
%   v. 79–80   Turn to the state of letters: why the ancients, ignorant of
%              nature's laws, reached perfection in every genre of writing —
%              the sense of beauty is drawn from the laws of man's own
%              intellectual nature; genius and taste.
%
% Numbering on the leaves, as read. (1) The Bibliothèque's foliation, thin
% grey figures at the top right of each recto, one per leaf: 30 (v. 61),
% 31 (v. 63), 32 (v. 65), 33 (v. 67), 34 (v. 69), 36 (v. 73), 37 (v. 75),
% 38 (v. 77), 39 (v. 79), given as \folio{}; at v. 71 the second figure is
% ambiguous, so no \folio there, only a note. Versos carry none. (2) At the top left of every view,
% v. 66 included, a struck number of two digits; mostly illegible under the
% stroke, but 60 (v. 70), 66 (v. 76) and 69 (v. 79) can be made out, and the
% others are consistent with a running page series one per view; the series is
% described, not transcribed, and no number was filled in from the pattern.
%
% Marks not in Germain's hand. At four places a name in a large hand,
% underlined, stands in the text at a paragraph end or head: « Varlé » (v. 61),
% « Gutmann » (v. 69), « Deraché » (v. 73), « Pemeguin » (v. 79), all
% uncertain readings; with them a boxed « Fl. 7 » in the margin (v. 69), a
% faint pencilled « Ch. 5 … » (v. 79), square brackets in the margin before
% a paragraph (v. 61, 70, 80), and an « avait » in upright letters (v. 70).
% They look like printing-house marks (compositors' names at the start of
% each take) — this is the pass's inference, not checked against anything.
% They are recorded in \note{}s where they fall and not transcribed as text.
%
% Notation: none; « --- » renders the long dashes she draws between
% paragraphs and at line ends. Her spelling is kept (« connoissance »,
% « tems », « hazardé », « satisfesantes », « apperçoit », « entr'eux »).
%
% The hand: a regular, rapid cursive, well inked; the struck words are
% usually legible under a single stroke. Not read: struck words at v. 61,
% v. 63 (two), v. 68 and v. 69; a
% word under an ink blot after « subit » (v. 73); a number overwritten at
% v. 72 (« plus de … sur cent »); a struck word before « premiers ouvrages »
% (v. 80); marginal struck words at v. 78; an illegible word before « de
% repos » (v. 68).
%
% Nothing here was checked against any published transcription.
\documentclass[11pt,a4paper]{article}
% The preamble shared by every transcription of the mathematicians' manuscripts in Gallica.
%
% It rests on one idea: **uncertainty compounds, it does not resolve.** A
% transcription that smooths over a doubtful word has destroyed the only thing
% it was made for — a reader must be able to tell, at every point, what was on
% the page and what was supplied. The apparatus macros below are therefore the
% critical apparatus, not decoration.
%
% The same commands are recognised by `scripts/render.mjs`, which produces the
% site's reading view, and by `scripts/tei.mjs`. A macro added here must be
% added there too, or rendering fails loudly — which is the wanted behaviour.
%
% Comments here are English, like the rest of the repository. The documents
% this preamble sets are French, and that is deliberate: see the skills.

\ProvidesPackage{germain}[2026/09/15 Transcriptions of mathematicians' manuscripts in Gallica]

\usepackage{amsmath,amssymb,amsthm}
\usepackage{mathrsfs}
% Kept from the parent preamble so the renderer's diagram support stays
% available; these manuscripts draw no commutative diagrams, and nothing here uses it.
\usepackage{tikz-cd}
\usepackage[english,french]{babel}
\usepackage{fontspec}

% --- Greek ------------------------------------------------------------------
% The text face stays fontspec's default, Latin Modern, which has no Greek:
% XeTeX drops every glyph it lacks with a line in the log and nothing on the
% page, and the Greek words the manuscripts quote vanished from the PDFs. So
% the Greek and Greek Extended blocks get a XeTeX character class of their own,
% and entering or leaving a run of it switches the family to CMU Serif — the
% same Computer Modern design, with polytonic Greek — and back. Only the
% family changes: a Greek word in \textit or \textbf keeps its shape and series.
% The transitions to and from every other class carry the tokens that class
% has towards an ordinary letter, so French spacing before « ; » or inside
% « » still holds after a Greek word. No grouping, so no brace or \endgroup
% can come out unbalanced; maths is left alone, since XeTeX inserts nothing
% there. Kept identical in `exercices.sty`.
\makeatletter
\newfontfamily\germainGreekfont{cmun}[
  Extension=.otf, NFSSFamily=germaingreek,
  UprightFont=*rm, ItalicFont=*ti, SlantedFont=*sl,
  BoldFont=*bx, BoldItalicFont=*bi, BoldSlantedFont=*bl]
\newXeTeXintercharclass\germain@greek
\def\germain@greekrange#1#2{%
  \@tempcnta=#1\relax
  \loop
    \XeTeXcharclass\@tempcnta=\germain@greek
    \ifnum\@tempcnta<#2\advance\@tempcnta\@ne\repeat}
\germain@greekrange{"0370}{"03FF}
\germain@greekrange{"1F00}{"1FFF}
\def\germain@greekprev{\rmdefault}
\def\germain@greekin{%
  \edef\germain@greekprev{\f@family}\fontfamily{germaingreek}\selectfont}
\def\germain@greekout{\fontfamily{\germain@greekprev}\selectfont}
\def\germain@greekjoin{%
  \XeTeXinterchartokenstate=\@ne
  \@tempcnta=\z@
  \loop
    \ifnum\@tempcnta=\germain@greek\else
      \XeTeXinterchartoks\germain@greek\@tempcnta=\expandafter{%
        \expandafter\germain@greekout\the\XeTeXinterchartoks\z@\@tempcnta}%
      \XeTeXinterchartoks\@tempcnta\germain@greek=\expandafter{%
        \the\XeTeXinterchartoks\@tempcnta\z@\germain@greekin}%
    \fi
    \ifnum\@tempcnta<4095\advance\@tempcnta\@ne\repeat}
% babel-french sets its own classes, and saves and restores the interchar
% state, each time a language is selected: join again after it does.
\addto\extrasfrench{\germain@greekjoin}
\addto\extrasenglish{\germain@greekjoin}
\AtBeginDocument{\germain@greekjoin}
\makeatother

\usepackage{geometry}
\geometry{a4paper,margin=25mm,marginparwidth=28mm}
\usepackage{marginnote}
\usepackage{xcolor}
\usepackage[normalem]{ulem}
\setlength{\emergencystretch}{3em}
\usepackage{hyperref}
\hypersetup{colorlinks=true,linkcolor=black,urlcolor=black,pdfborder={0 0 0}}

% --- Batch metadata ---------------------------------------------------------
% Taken as they stand by the rendering script, which shows them at the head of
% the reading view. They fix what the file claims to cover. The macro names are
% the parent project's, so that the scripts stay shared; \folder here means the
% volume's slug — `fr-9115` — and \pages counts Gallica views.

\newcommand{\folder}[1]{\gdef\grothFolder{#1}}
\newcommand{\batch}[1]{\gdef\grothBatch{#1}}
\newcommand{\pages}[2]{\gdef\grothFirst{#1}\gdef\grothLast{#2}}
\newcommand{\foldertitle}[1]{\gdef\grothFoldertitle{#1}}

% The holder's own shelfmark, printed as they write it: « Français 9115 ».
\newcommand{\shelfmark}[1]{\gdef\grothShelfmark{#1}}

% The dating, copied verbatim from the holder's catalogue — never inferred
% here. « XIXe siècle » is the BnF's; a pass that dates a leaf from its content
% says so in a \note{}, as its own inference.
\newcommand{\dating}[1]{\gdef\grothDating{#1}}

% The status notice, printed in the title block and repeated as a diagonal
% watermark on every page of the PDF. The manuscripts are in the public
% domain; what this line declares is not a right but a quality — an
% unverified first pass by a machine. The skills write the line; a file
% without it gets no watermark, so leaving it out is visible in the output.
\newcommand{\watermark}[1]{\gdef\grothWatermark{#1}}

\gdef\grothFolder{?}\gdef\grothBatch{}\gdef\grothFirst{?}\gdef\grothLast{?}%
\gdef\grothFoldertitle{}\gdef\grothShelfmark{}\gdef\grothDating{}\gdef\grothWatermark{}

% --- The résumé -------------------------------------------------------------
\newenvironment{resume}
  {\small\setlength{\parskip}{0.45em}}
  {\par\medskip\hrule\medskip}

% --- Keywords ---------------------------------------------------------------
% In English, closing the modernised reading's résumé: the modern vocabulary
% under which to search for what the leaves construct. The manifest extracts
% them to tag the volume — this line is the single source of the tags.
\newcommand{\keywords}[1]{\par\smallskip\noindent
  {\footnotesize\textsc{Keywords} --- \textit{#1}}\par}

% --- The critical apparatus -------------------------------------------------

% The Gallica view number, in the margin. This is the anchor that lets a
% disagreement over a formula be settled by looking at *one* image rather than
% a volume of 750. The site uses it to turn the facsimile as the transcription
% is scrolled. A view is one scanned image: usually one side of a leaf.
\newcommand{\page}[1]{%
  \marginnote{\footnotesize\color{gray!70}\textsf{v.\,#1}}%
  \label{page:#1}\ignorespaces}

% The BnF's pencilled foliation, where the leaf carries it and a pass can read
% it: « 198r ». This is what the literature cites — « ff. 198r–208v » — and
% the one line that turns a citation into a view. Recorded, never inferred:
% a folio a pass cannot read is not written.
\newcommand{\folio}[1]{%
  \marginnote{\footnotesize\color{gray!50}\textsf{f.\,#1}}\ignorespaces}

% The modernised reading's coarser counterpart to \page{}: which views a
% section is drawn from.
\newcommand{\pagerange}[2]{%
  \marginnote{\footnotesize\color{gray!70}\textsf{v.\,#1--#2}}%
  \ignorespaces}

% Illegible: never guessed.
\newcommand{\ill}{\textcolor{red!60!black}{[\,\ldots\,]}}

% A doubtful reading: the word is given, and flagged as uncertain.
\newcommand{\uncertain}[1]{\uline{#1}}

% An editorial addition — a missing letter, an implied word.
\newcommand{\add}[1]{\textcolor{blue!50!black}{[#1]}}

% Struck out by the writer. What was crossed out is part of the document.
\newcommand{\struck}[1]{\sout{#1}}

% The transcriber's note, on the state of the page or the mathematics.
\newcommand{\note}[1]{\par\smallskip{\small\color{gray!60!black}\textit{[#1]}}\par\smallskip}

% A marginal note of hers, distinct from the body of the text.
\newcommand{\marginal}[1]{\par\smallskip\noindent\hspace{1em}%
  {\small\color{gray!60!black}#1}\par\smallskip}

% --- Title ------------------------------------------------------------------
% The title block is French, like the documents themselves.

\AddToHook{shipout/background}{%
  \ifx\grothWatermark\empty\else
    \begin{tikzpicture}[remember picture, overlay]
      \node[rotate=52, scale=3.4, align=center, text=gray!18,
            font=\sffamily\bfseries]
        at (current page.center) {\grothWatermark};
    \end{tikzpicture}%
  \fi}

\AtBeginDocument{%
  \begin{center}
    {\large\bfseries Manuscrits de mathématiciens (Gallica) ---
      \ifx\grothShelfmark\empty\grothFolder\else\grothShelfmark\fi}\\[2pt]
    {\itshape\grothFoldertitle}\\[4pt]
    {\small vues \grothFirst--\grothLast
      \ifx\grothBatch\empty\else\ (lot \grothBatch)\fi}\\[2pt]
    \ifx\grothDating\empty\else
      {\small\color{gray!60!black}Datation du catalogue : \grothDating}\\[2pt]
    \fi
    \ifx\grothWatermark\empty\else
      {\small\bfseries\color{red!45!gray}\grothWatermark}\\[2pt]
    \fi
    {\footnotesize\color{gray!60!black}Transcription établie d'après les images IIIF de Gallica.
     Source gallica.bnf.fr / Bibliothèque nationale de France.\\
     Les lectures incertaines sont soulignées, les illisibles notées [\,\ldots\,],
     les ajouts entre crochets ; \textsf{v.} numérote les vues de Gallica,
     \textsf{f.} la foliotation au crayon de la BnF.}
  \end{center}
  \vspace{1em}}

\endinput


\folder{fr-9117}
\shelfmark{Français 9117}
\batch{4}
\pages{61}{80}
\dating{1801-1900}
\watermark{Lecture automatique\\non vérifiée}
\foldertitle{« Considérations générales sur l'état des sciences et des lettres aux différentes époques de leur culture, » par Sophie GERMAIN.}

\begin{document}

\note{Numérotation des feuillets. La foliotation de la Bibliothèque, en haut
à droite de chaque recto, donne 30, 31, 32, 33, 34, 36, 37, 38, 39 aux vues
61, 63, 65, 67, 69, 73, 75, 77, 79 ; à la vue 71 le second chiffre est ambigu
et aucun folio n'est porté. En haut à
gauche de chaque vue, y compris la vue 66, un nombre de deux chiffres biffé,
le plus souvent illisible sous la biffure ; seuls 60 (vue 70), 66 (vue 76) et
69 (vue 79) se lisent. La vue 66, verso blanc du feuillet écrit à la vue 65, n'est
pas transcrite.}

\page{61}\folio{30r}
\note{En tête de la page, au-dessus de
la première ligne, un mot souligné, lu \uncertain{Varlé}, dont le rôle n'est
pas établi ; à gauche, un signe en forme de « § » et une grande accolade qui
embrasse les deux premières lignes.}
qui a séduit autrefois l'esprit humain.

De toutes nos idées, la plus abstraite est celle de l'être ; car, celle
du néant est toute négative. L'être nous appartient ; il pénètre notre
intelligence et l'éclaire du flambeau de la vérité. Les idées du beau, du
bon, sont plus compliquées ; nous les devons à la comparaison entre les
connoissances acquises et notre modèle intérieur. D'autres idées sont plus
immédiatement dues à nos sensations. Ainsi, le grand, le petit, le fort, le
foible, expriment des comparaisons qu'il seroit absurde de regarder comme
innées. J'en dirais autant de ce que nous appelons beauté ou bonté
relatives ; ces connoissances sont toutes acquises à l'aide des sensations
et de la réflexion.

C'est à l'uniformité des conditions de l'être qu'il faut
attribuer\note{« attribuer » est écrit dans la marge, en tête de la ligne,
suivi de quelques lettres biffées.} \struck{\ill{}} le sentiment d'analogie
qui dirige toutes les opérations de notre entendement. L'histoire de
l'esprit humain nous apprend que ce sentiment a produit des erreurs
grossières\note{« grossières » ajouté dans l'interligne.} aussi bien que
\struck{les plus} d'heureuses pensées.

On peut demander comment une cause dont l'action est constante a cependant
amené des résultats si différens ?

Nous allons voir que de telles différences \struck{étoient} sont le
résultat inévitable des manières diverses dont on a cherché à réaliser les
indications vers lesquelles nos tendances intellectuelles n'ont jamais cessé
de nous entraîner.

Par leur nature les conditions de l'être sont abstraites : et, s'il en
étoit autrement on ne concevroit pas qu'elles fussent universelles. L'esprit
de système consistoit à prendre un fait connu, c'est-à-dire une vérité
particulière, pour base d'un ordre de

\page{62}
faits. Ceux-ci, on ne les considéroit plus en eux-mêmes ; on y cherchoit
seulement \struck{ou} les rapports vrais ou supposés qui les lioient au
premier ; ainsi, en assemblant un certain nombre d'êtres particuliers, on
attribuoit à l'un d'eux la domination des autres ; en sorte que ces
derniers \struck{étoient} dépouillés de leurs réalités individuelles,
\struck{et revêtus} étoient revêtus de celle qui convenoit uniquement à la
vérité dominante dont on avoit fait\note{« fait » ajouté dans l'interligne.}
choix.

\struck{Ainsi} Au lieu de chercher des analogies, on vouloit trouver des
identités, parce qu'en effet des identités seroient plus simples que des
analogies, et, par conséquent \struck{elles seroient} plus
satisfesantes.\note{Un trait courbe, tiré sous « que des analogies, et, par
conséquent », rejoint « plus satisfesantes » à la ligne suivante.} Le type
\uncertain{du} vrai, l'unité de l'être, l'ordre, la proportion des parties
dont la nécessité s'est toujours fait sentir, on croyoit pouvoir les
réaliser arbitrairement, au gré d'une imagination capricieuse.

On devoit s'égarer ; et pourtant les erreurs de l'esprit humain, qui
sembleroient devoir être inépuisables, se sont toutes rapprochées de
certaines vérités, et n'ont pas été aussi nombreuses que le vice des
procédés pourroit le faire croire. C'est qu'en effet le sentiment du vrai
n'a jamais abandonné les auteurs de tous ces systèmes. Cet heureux
sentiment n'a pas suffi pour les préserver des suppositions arbitraires et
forcées ; mais \struck{pourtant} il a servi à \struck{mettre des bornes}
retenir leur imagination dans de certaines limites.

À l'esprit de système succèdent aujourd'hui \struck{le l'esprit méthodique}
les recherches méthodiques : la généralité des conditions de l'être est
mieux comprise, dans chaque sujet ; on dirige ses efforts vers la découverte
de leur

\page{63}\folio{31r}
réalisation ; mais on ne les confond plus avec \struck{l'existence} les
conditions particulières qui appartiennent en propre à\note{« les
conditions » et « qui appartiennent en propre à » ajoutés dans l'interligne ;
un « de » biffé devant « la vérité ».} la vérité individuelle \struck{qui a
d'abord servi à} dont la découverte fortuite a décelé l'existence d'un ordre
de phénomènes longtems inaperçu.

L'expérience est consultée ; on cherche d'abord à multiplier les faits en
variant les circonstances dans lesquelles ils peuvent se manifester.
\struck{Avertis} par ce sentiment intime de l'analogie de l'existence des
lois qu'\struck{on} il n'apperçoit pas encore, \struck{il}
l'observateur\note{« averti » et « l'observateur » sont ajoutés dans
l'interligne, au-dessus du début de la ligne ; la phrase n'a pas été remise
en ordre, et la place exacte de « averti » n'est pas marquée.} s'applique à
séparer les circonstances qui compliquent les résultats, en cherchant pour
chacune d'elles les plus grandes et les moins grandes influences. Alors les
faits se classent, ils présentent un enchaînement, un ordre \struck{et
enfin} les lois \struck{se mani} dont l'existence avoit été
\uncertain{prévues} se manifestent, et une branche nouvelle de la science
vient s'ajouter à des connoissances plus anciennes.

À cette période\note{« À cette période » ajouté dans l'interligne, au-dessus
de « On ».} on ne possède cependant encore\note{« encore » biffé en fin de
ligne et récrit au-dessus.} que la partie \struck{\ill{}} expérimentale. La
théorie est créée lorsque, la nature des faits \struck{s'étant} prêtée à une
expression analytique, on est parvenu à tirer de cette expression des
conséquences conformes à l'expérience\note{« à l'expérience » ajouté dans
l'interligne.} : \struck{en même tems qu'elles sont explicatives, des faits
observés.} \struck{Alors} et les formules nées des premières observations
\struck{viennent} révèlent ensuite l'existence de faits encore inconnus.

Aujourd'hui \struck{\ill{}} que diverses branches de la physique sont
entrées dans le domaine des sciences mathématiques,\note{Addition écrite
dans l'interligne et poursuivie dans la marge de gauche (« le domaine des
sciences mathématiques »), où « aujourd'hui » est encadré ; un trait la
rattache à la ligne.} on voit avec admiration, les mêmes intégrales,
\struck{différentes} à l'aide des constantes fournies par divers genres de
phénomènes, représenter des faits entre lesquels on n'auroit jamais supposé
la moindre analogie. Leur ressemblance devient alors sensible ; elle est
intellectuelle, elle dérive des lois de l'être ; et \struck{ce} ce qui fut
autrefois le rêve d'une imagination hardie ; incertaine encore des formes
qu'elle devoit revêtir, l'identité des rapports de l'ordre et des
proportions dans des existences les plus diverses, vient se montrer aux yeux
\struck{avec}

\page{64}
en même tems qu'à la pensée avec l'évidence qui appartient aux sciences
exactes.

Mais les lois de l'être ne régissent pas seulement les faits qui sont du
domaine des sciences ; elles s'appliquent également à l'ordre intellectuel.
--- C'est en s'approchant de plus en plus du type de l'être ou du vrai,
source de toutes nos connoissances réelles, que les théories se
perfectionnent ; que la morale s'épure ; que la politique s'éclaire ; que la
métaphysique cesse de s'égarer ; que la littérature et les beaux arts se
rendent compte des règles qu'ils ont pratiquées et des grands effets qu'ils
peuvent produire.

Malgré l'extrême différence des genres, toutes ces choses ont entr'elles
des rapports d'ordre et de proportion qui deviennent d'autant plus sensibles
qu'elles sont examinées de plus près.

Si, par des progrès, qui semblent encore aujourd'hui au delà de toutes
espérances raisonnables, la langue des calculs devenoit applicable à des
questions morales, politiques, métaphysiques, ou à celles qui tenant de plus
près à notre manière de sentir, \struck{forment} \struck{le} composent le
domaine du goût, la ressemblance des formules rendroit évident que des
objets si divers, ont entr'eux la ressemblance que \struck{leur} impriment
les lois de l'être. Leur nature spéciale seroit représentée par des
constantes ; \struck{et} toutes \struck{ce qui concerne} les propositions
relatives à chaque sujet se trouveroient exprimées par des fonctions dont
les formes se reproduiroient sans

\page{65}\folio{32r}
cesse et offriroient, par leur identité, la preuve complète des
ressemblances intellectuelles dont nous parlons.

\struck{Pr} Choisissons un exemple qui puisse faire mieux comprendre notre
proposition.

Nous voyons, dans différens genres de phénomènes, que la tendance à la
régularité se manifeste \struck{dans} par les formules qui leur sont
applicables ; car les termes qui \struck{renferment} expriment
l'irrégularité \struck{sont} renferment la \struck{tems} durée, de manière
à montrer, qu'après un tems fort court ils doivent disparoître.
\uncertain{Si bien} le théorème relatif à la courte durée de l'action des
causes perturbatrices seroit attesté, dans notre supposition, par les formes
du calcul.

On verroit en morale combien peu doivent durer les effets de la fraude, du
mensonge et de l'injustice. Il deviendroit sensible que le vrai et le juste
tendent sans cesse à faire disparoître les obstacles qui s'opposent à leur
manifestation.

En politique, parmi les causes qui agissent sur le système on
distingueroit\note{« on distingueroit » ajouté dans l'interligne.} quelles
sont celles qui, étant dues à des forces toujours croissantes, doivent finir
par prédominer ; tandis que des causes accidentelles, dont l'effet est fort
grand à un instant donné, doivent au contraire cesser entièrement leur
action après un tems fort court.

\note{Un crochet dans la marge embrasse les deux lignes qui suivent.}
Dans les sciences de raisonnement on trouveroit également que l'erreur tend
à disparoître.

En matière de goût, la mode est une cause perturbatrice ; aussi son empire
n'est-il pas de longue durée.

Il est donc vrai que, quelques divers que soient les sujets, les actions qui
troublent l'ordre naturel tendent à s'anéantir.\note{« tendent à
s'anéantir » écrit sous la dernière ligne, faute de place.}

\page{67}\folio{33r}
L'analogie qui se fait remarquer \struck{dans} entre les différens objets
dont nous avons connoissance ne se borne pas à un seul point. On pourroit
affirmer, par exemple, que la mécanique rationnelle toute entière offre avec
les sciences politiques des ressemblances telles que les théorèmes dont se
compose la première, \struck{offre} deviennent, par rapport aux secondes des
propositions dont la vérité est incontestable.

Ainsi\note{« Ainsi » ajouté dans l'interligne, précédé d'un crochet.}
\struck{le repos n'ait} l'équilibre entre plusieurs forces vient de ce que
l'action des unes est opposée de direction et égale en puissance à celle des
autres ;\note{« à celle des autres » ajouté dans l'interligne.}
\struck{elles} se compensent et se décomposent ; elles produisent alors des
résistances dans un sens qui n'est pas celui de leur action directe.

La même chose a lieu à l'égard des forces qui \struck{composent} naissent de
l'état de société. Si elles sont opposées de direction et égales en
puissance, l'état de repos se maintient de lui-même. --- Il y a de l'art à
changer le sens dans lequel elles agissent, \struck{par} en leur
opposant\note{« en leur opposant » ajouté dans l'interligne.} des obstacles
indirects. Le parallélogramme des forces pourroit servir d'emblême à ce
genre d'adresse.

Lorsqu'un système est en repos, cet état peut être dû à des conditions
essentiellement différentes. Si une cause extérieure vient à agir sur le
système, ou il tendra à reprendre sa position initiale, et
\struck{reviendra} l'équilibre se rétablira au moyen d'oscillations, dont
l'amplitude diminuera à chaque instant ; ou bien le mouvement

\page{68}
communiqué éloignera de plus en plus le système de sa position initiale, et
\struck{il} le système ne reviendra à l'état \ill{} de repos \struck{que
lorsqu'il sera} qu'après être\note{« en plus », « le système » et « qu'après
être » ajoutés dans l'interligne.} parvenu à une situation entièrement
différente. ---

Les deux cas d'équilibre stable et d'équilibre non stable se font également
remarquer dans l'état social ; on voit des causes propres à l'agiter
produire tantôt de légers mouvemens, qui cessent d'eux-mêmes ; tantôt des
révolutions complètes, \struck{en sorte que} qui ne permettent à l'état de
paix intérieure de renaître \struck{n'arrive plus} qu'après de grands
changemens dans l'ordre social.\note{« qui ne permettent à » et « de
renaître » ajoutés dans l'interligne.}

Si on veut pousser la comparaison plus loin, l'analogie ne se démentira pas.

L'équilibre stable a lieu lorsque tous les points du système ont atteint la
situation qui convient à leur tendance naturelle. ---

L'état de tranquillité est durable lorsque tous les individus qui composent
la société sont dans la situation qui convient à leur tendance naturelle.

L'équilibre est non stable \struck{lorsqu'il} quand il est établi sur un
point où il ne peut subsister qu'autant de tems qu'il est à l'abri de tout
choc ; \struck{en sorte} et \struck{\ill{}} que \struck{que} le
\struck{mouv} moindre \struck{dévelop} dérangement venant à rendre aux
points qui le composent la liberté de se mouvoir dans la direction de leur
tendance naturelle, l'état initial doit finir par être changé en un état
opposé au premier ; \uncertain{en} \struck{continuer} sorte que le mouvement
ne \struck{peut} puisse cesser avant que ce nouvel état, qui n'est autre que
celui \struck{auquel appartient} qui constitue l'équilibre stable,
\struck{est} ait remplacé l'état initial. ---\note{« sorte », « puisse »,
« qui constitue » et « ait » ajoutés dans l'interligne ; un astérisque suit
« qui constitue », sans renvoi visible sur la page.}

\page{69}\folio{34r}
Les États gouvernés sans égard aux tendances sociales conservent la
tranquillité intérieure aussi longtems qu'aucun événement ne vient agiter les
esprits, mais la moindre circonstance suffit pour \struck{agiter} ébranler la
société jusque dans ses fondemens. Nous voyons alors chaque volonté
individuelle recevoir une impulsion nouvelle, et les mouvemens qui en sont
la suite subsister jusqu'à ce que l'état reconstitué sur des bases
\struck{\ill{}} plus solides offre à chacun les garanties dont il avoit senti
le besoin.

\note{À la fin de ce paragraphe, à droite, un nom souligné, d'une écriture
plus grande, lu \uncertain{Gutmann} ; dans la marge, devant le paragraphe
suivant, une mention encadrée, lue \uncertain{Fl. 7}, et le premier mot
« Dans » encadré aussi. Ces marques, comme le mot souligné en tête de la
vue 61, ne font pas partie du texte ; leur rôle n'est pas établi ici.}
Dans un système de points doués de pesanteur, chacun tend à se placer aussi
près que possible du centre de la terre. La situation qu'ils atteignent
n'est pas celle qu'ils obtiendroient s'ils étoient libres ; mais elle dépend
à la fois de leur liaison et de leur tendance individuelle. ---

Dans l'état social chaque individu tend vers le bien être et la première
condition à remplir est que le bien être de chacun nuise le moins possible à
celui des autres.

L'équilibre d'un système exige que le centre de gravité soit appuyé. S'il se
trouve placé le plus bas possible, l'équilibre est stable. ---

Le repos d'un état seroit impossible à maintenir si on n'avoit aucun égard à
la tendance de l'époque, ou, ce qui est la même chose, à l'opinion. Il faut,
ou lui opposer de puissans obstacles, ou savoir se conformer à ses
exigences. Ces deux manières d'envisager la question conduisent à la
tranquillité précaire ou à la tranquillité durable.

\page{70}
Considérons présentement les effets de l'impulsion.

Si la direction du\note{« direction du » ajouté dans l'interligne, sur
« le » biffé.} mouvement communiqué à un système de corps passe par le
centre de gravité du système, il sera mû comme si tous les points qui le
composent étoient réunis en un seul, et la force toute entière sera employée
à produire l'effet qu'on en attendoit. ---

--- De même aussi lorsque l'action du gouvernement est dirigée dans le sens
de l'opinion, la société semble se mouvoir comme un seul individu, qui
agiroit conformément à ses intérêts ; et les forces de l'état sont employées
toutes entières au profit de la prospérité générale.

S'il arrivoit que la direction du mouvement fût différente, la force seroit
décomposée en deux portions. L'une, celle dont la direction passeroit par le
centre de gravité du système, agiroit comme si elle étoit seule pour faire
avancer le système dans la route où l'on auroit voulu le pousser ; tandis que
l'autre portion de la force motrice, entièrement perdue par rapport à ce but
n'auroit d'autre effet que de faire tourner le système autour de son centre
de gravité\ldots{} \struck{Enfin s'il} si \struck{arrivoit que} l'impulsion
\struck{eut} avait\note{« avait » est écrit dans l'interligne en lettres
droites, d'une écriture qui ne paraît pas celle du texte.} été assez
maladroite pour que la première portion de la force motrice fût nulle, le
système n'auroit aucun mouvement progressif ; la force de rotation
subsisteroit seule, et il seroit dans la nature de cette force
\struck{d'éloigner les} de tendre à détruire la liaison entre les

\page{71}
\note{Foliotation de la Bibliothèque, en haut à droite : « 3 » suivi d'un second
chiffre ambigu, qu'on attendrait 5 dans la série mais qui ne se lit pas
sûrement ; aucun folio n'est donc porté ici.}
les diverses parties du système. ---

--- Nous voyons aussi l'action des gouvernemens être en partie favorable et
en partie nuisible lorsque ses \struck{effort} efforts \struck{obéissent}
satisfont en quelques points à l'opinion publique, qu'ils contrarient sous
d'autres rapports. ---

--- S'il pouvoit exister une administration assez malavisée pour agir en
toutes circonstances dans des directions opposées à l'opinion, ou ce qui est
la même chose, à l'intérêt public, l'État \struck{confié aux soins}
éprouveroit une agitation intérieure qui tendroit à le dissoudre. Ainsi, par
exemple, il se pourroit, qu'à la première occasion les provinces frontières
favorisassent les prétentions de l'état voisin qui voudroit les envahir ;
car, en politique, aussi bien qu'en mécanique les points de limites sont
ceux qui se trouvent le plus agités dans les mouvemens dont nous parlons. Les
forces tangentielles sont nulles au centre du système ; le désir de la
séparation \struck{est} seroit absurde dans les capitales.

Les sociétés sont composées de trois élémens principaux. On y trouve des
intérêts, des passions, de l'inertie. Les individus \struck{rassemblent}
réunissent quelquefois les trois manières d'être correspondantes ; mais l'une
d'elle domine le plus souvent ; et elles\note{« elles » ajouté dans
l'interligne.} forment autant de caractères différens. Ces trois caractères
présentent des ressemblances avec la manière dont se comportent les corps
durs, élastiques et mous. Ainsi on voit les personnes exclusivement occupées
de leurs intérêts tenir obstinément à la route qui les conduit au but
qu'elles veulent atteindre, et résister à tout mouvement contraire ; en sorte

\page{72}
que \struck{les} l'obstacle qu'elles rencontrent ne les détermine à changer
de direction que lorsqu'ils \struck{ont} a détruit toute leur force. Au
contraire, les personnes mues par leurs passions prennent au moindre obstacle
un parti inattendu ; elles changent de route, et agissent d'une manière toute
contraire à celles qu'elles avoient d'abord adoptées. Enfin nous voyons
certains caractères, amis du repos souffrir des lésions réelles, plutôt que
de songer à réagir.

Dans les tems de tranquillité, les intérêts dominent. L'administration doit
tendre à les protéger ; ce soin semble ne devoir pas offrir de grandes
difficultés ; car il est dans la nature des intérêts, d'indiquer eux-mêmes
les mesures qui leur sont favorables ; leur direction est connue et
invariable ; ils servent de base à l'opinion publique.

\struck{Au contraire, lorsque} Mais, que\note{« Mais, que » ajouté dans la
marge de gauche.} le repos intérieur de l'état \struck{vienne à être} soit
troublé, les passions, aisément maintenues dans l'état de paix, viennent
augmenter le trouble ; elles agissent dans mille directions à la fois ; on
ne sait où elles tendent ; et il est \struck{alors} fort difficile de prévoir
quel sera le résultat de leur choc.

On n'a pas encore imaginé de faire une statistique des caractères. Mais on
peut être sûr qu'il y a \struck{sur cent} un assez grand nombre d'hommes qui
agissent toujours conformément à leurs intérêts ; plus de \ill{} sur cent
peut-être\note{Le nombre, écrit en chiffres et surchargé, n'a pas pu être
lu ; « sur cent » est récrit dans l'interligne au-dessus, après avoir été
biffé deux lignes plus haut. Un long trait termine la ligne.}

\page{73}\folio{36r}
L'autre partie est partagée en deux portions. L'une \struck{composée} des
êtres irritables pour qui les intérêts semblent toujours méprisables
\struck{lorsqu'ils sont en} comparés avec l'objet de leurs passions. Suivant
les âges et les positions, \struck{elles} ces passions\note{« ces passions »
ajouté dans l'interligne.} peuvent prendre des caractères différens ; mais
l'amour propre est la plus constante de toutes. \struck{Une} L'autre
portion est \struck{composée} formée des personnes qui se rend\struck{ent}
aisément esclaves de leurs habitudes, craignent tout ce qui les pourroit
changer ; \struck{elles ne} ceux-ci connoissent ni l'ambition des richesses
ni celle de la gloire, ni les affections vives ; ce sont des gens inertes.

\note{À la fin de ce paragraphe, un nom d'une grande écriture, souligné
d'un paraphe et taché d'encre, lu \uncertain{Deraché}.}
Mais \struck{ainsi qu'}il n'existe dans la nature morte aucun corps qui
\struck{soit} soit parfaitement dur, c'est-à-dire qui ne puisse changer de
formes sous des efforts puissans et répétés ; aucun corps qui soit
parfaitement élastique, c'est-à-dire qui ne retienne rien de la direction
dans laquelle on le pousse ; aucun corps parfaitement mou, c'est-à-dire
aucun corps que le choc ne puisse faire changer de \struck{forme} place, et
qui absorbe toute la force employée par le seul changement de forme qu'il
subit\ill{}.\note{Une tache d'encre couvre un mot après « subit ».} De
même\note{« De même » ajouté dans l'interligne.} on ne voit pas non plus de
gens tellement attachés à l'intérêt qu'en certain moment de leur vie, ils ne
puissent agir par d'autres motifs : \struck{et de même} les hommes
passionnés cèdent \struck{dans certains} quelquefois à leurs intérêts : et
les personnes naturellement amies du repos peuvent trouver dans les choses
et dans les personnes qui les environnent matière à exciter en
elles\note{« en elles » dans la marge de gauche.} le désir de la richesse,
celui de la renommée ou de l'affection ; \struck{leurs} passions de ces
derniers\note{« de ces derniers » ajouté dans l'interligne.} seront
faibles ; mais

\page{74}
enfin elles peuvent n'être pas absolument sans effets extérieurs.

Si bien le trouble d'un état rompt la balance habituelle entre les trois
nombres qui représentent ces caractères différens. Tous les individus
reçoivent une impulsion, qui les transforme en gens passionnés. Le mouvement
se distribue sans doute inégalement entr'eux ; mais il est extrêmement
difficile d'évaluer la force des masses composées d'élémens nouveaux. Les
directions sont incertaines et variables. L'agitation se manifeste surtout
dans des actions instantanées ; et cette circonstance \struck{double}
redouble la difficulté. En effet, dans aux époques de paix et de tranquillité
publique, ceux qui tiennent les rênes ont tout le tems nécessaire pour
résoudre les mesures qu'ils doivent prendre. Des lumières, de l'habitude des
affaires, l'intention de faire le plus de bien avec le moins de mal
possible, suffisent pour gouverner avec habileté. Dans les momens de crise,
c'est toute autre chose. Les circonstances deviennent pressantes ; il faut
savoir se résoudre promptement. On a souvent aussi besoin de courage ; et le
courage ne se trouve pas nécessairement joint aux qualités qui
\struck{dans tout autre} font l'homme habile. La société court donc mille
dangers, qu'il est aussi difficile d'éviter que de prévoir.

\page{75}\folio{37r}
Ajoutons \struck{à tout cela} que des individus doués de grandes forces par
la nature se trouvoient placés durant le calme dans des positions qui
annuloient ces forces ; tandis qu'à la faveur du trouble, ils apparoissent de
tous côtés, armés d'une énergie jusque là inconnue. De tels individus
n'avoient pas prévu qu'ils sortiroient un jour de la nullité à laquelle leur
position sociale les avoit condamnés ; ils ne se sont livrés à aucune étude
spéciale avant de prendre place parmi les hommes qui influeront sur le sort
de leurs semblables ; les partis violens sont les seuls qu'ils puissent
adopter, parce qu'ils y trouvent l'emploi de leurs forces, et
\struck{parce qu'ils sont} sont\note{« sont » ajouté dans l'interligne.}
dépourvus de l'adresse qui seroit le fruit des connoissances qui leur
manquent.

L'égalité est une erreur ; et la mécanique vient encore ici nous suggérer
une analogie nouvelle. Deux masses de même poids peuvent avoir des forces
vives très différentes ; le plus petit poids placé au bout d'un levier fera
équilibre à une masse aussi forte qu'on voudra ; \struck{il} ne s'agit que
d'établir l'égalité entre les forces virtuelles. ---

Nous voyons la même chose dans les sociétés ; et les révolutions ne sont si
dangereuses, si hazardeuses, si incertaines dans leur résultats que parce
qu'elles changent tout à coup

\page{76}
les rapports entre les forces vives des différentes classes de la société.
À la vérité, et nous l'avons dit plus haut, ce que les révolutions ont de
violent et d'irrégulier disparoît bientôt en vertu de ce théorème général
qui montre qu'en toute chose les forces perturbatrices sont fonctions du tems
et que la régularité tend à s'établir dans tout système, de quelque nature
qu'il puisse être.

Si nous voulons présentement jeter un coup d'\struck{œil} œil sur les
ressemblances qu'offrent entr'eux ce qu'on nomme l'ordre physique, moral, et
intellectuel, nous trouverons des analogies remarquables.

\struck{Les langues attestent ces analogies} Sans leur secours, le langage
figuré n'auroit pu naître.\note{« Sans leur » est écrit en bout de la ligne
biffée ; au-dessus de « secours », un mot biffé, lu \uncertain{desquelles},
remplacé par « de ces analogies », qui n'a pas été remis à sa place dans la
phrase.} Elles ont été senties dans tous les tems ; et peut-être même
n'auroit-on rien d'essentiel à ajouter aux remarques de ce genre que l'examen
des langues a suggéré. \struck{Cependant} Contentons nous d'observer combien
sont judicieuses ces applications du langage propre au langage figuré.

La force morale et la force intellectuelle se comportent en effet comme la
force physique. Il y a de part et d'autre des compositions et des
décompositions analogues, ainsi nous voyons nos diverses facultés, concourir
à un seul fait moral et intellectuel, comme il arrive à des forces

\page{77}\folio{38r}
de \struck{différentes} diverses natures et de directions différentes, de
donner une résultante qui, dans sa valeur et dans sa direction, représente la
véritable force motrice.

Mais cette justesse d'expressions paroît encore plus frappante lorsque,
fondée \struck{d'abord} sur un premier rapport apparent et connu, elle se
soutient à l'égard des rapports qu'on n'avoit pas eu \struck{en vue} d'abord
en vue. Prenons un exemple. Au physique, on appelle monstre un être
difforme ; cette expression, transportée au moral s'applique aux êtres
vicieux, parce que le vice est en effet une difformité morale. --- Lorsque
la langue s'est formée, l'anatomie n'existoit pas. Cette science nous a
appris que la monstruosité avoit pour cause le développement extraordinaire
de certains organes, qui, attirant à eux toutes les forces de la vie,
privoient les autres organes de la nourriture nécessaire au développement
qu'ils acquièrent dans l'état ordinaire.

En examinant de plus près les êtres qui effrayent les sociétés par de grands
attentats ou même ceux qui les troublent par des désordres habituels, nous
\struck{les} trouvons que des qualités hors de mesure les entraînent soit à
des forfaits, soit seulement à une \struck{contravention} infraction des lois
des sociétés. De telles qualités absorbent leur \struck{être}
moralité\note{« moralité » ajouté dans l'interligne, suivi de deux points.} :
ils manquent d'autres qualités qui, dans des individus moins généreux, moins
courageux, entretiennent l'amour de la justice

\page{78}
et celui de l'ordre.

On pourroit trouver d'autres exemples du même genre. Ils attestent cette
vérité fondamentale qu'un seul rapport bien constaté entre deux sujets de
genres différens en annonce un grand nombre d'autres.

Je ne sais si cette proposition énoncée formellement ne paroîtroit pas
hazardée ; mais chacun \struck{sait} raisonne\note{« raisonne » ajouté dans
l'interligne.} pourtant comme si elle étoit indubitable. \struck{Une sorte}
Elle renferme le principe de l'analogie ; principe qu'un de nos auteurs a
\struck{regardé comme} appellé une méthode d'invention, et qui, sans lui
refuser ce noble caractère, \struck{est employée est} peut être regardé comme
également propre aux emplois les plus communs, puisqu'\struck{il} elle
dérive du sentiment des lois de l'être, dont le caractère est partout
semblable.\note{« appellé » et « peut être regardé comme » ajoutés dans
l'interligne ; « une » biffé devant « méthode ». Les accords (« regardé »,
« il » / « elle ») sont restés tels que la correction les a laissés.}

L'habitude de l'étude nous donne une grande aptitude à saisir les
analogies ; et c'est en quoi elle nous sert à acquérir avec facilité des
connoissances nouvelles. Aux premiers mots sur\note{« sur » ajouté dans
l'interligne, au-dessus de « d'un » biffé.} un sujet encore inconnu, notre
esprit cherche à en fixer la nature, c'est-à-dire qu'il cherche à quel
module nouveau il va appliquer les lois qui \struck{leur} conviennent à
tout.\note{Dans la marge de gauche, en face de ces lignes, trois ou quatre
mots biffés et surchargés, illisibles, qui semblent une première rédaction
de l'addition.} Ce point étant fixé, notre esprit\note{« notre esprit »
ajouté dans l'interligne, au-dessus de « il » biffé.} avance à grands pas
\struck{dans} la route qui s'ouvre \struck{à ses yeux} devant lui. À chaque
instant, la règle de proportion trouve mille applications ; et si l'analogie
sert à l'invention, elle n'est pas moins utile à l'étude des sciences déjà
faites. Il n'appartient pourtant pas qu'on lui applique la dénomination de
\emph{méthode}. L'analogie n'est pas

\page{79}\folio{39r}
d'invention humaine ; elle existe par elle-même. Notre intelligence est
propre à la reconnoître ; elle aide nos premiers efforts ; elle instruit
l'enfant, quelquefois aussi elle l'induit en erreur et, quoiqu'il ne
s'agisse alors que des idées les plus communes, il est facile de voir que
les déviations de nos premiers jugemens, sont produites par la même cause qui
a enfanté \struck{fa} \struck{produit} les systèmes hazardés \struck{celle des
philosophes}.\note{« enfanté » est écrit dans la marge de gauche, « les
systèmes hazardés » dans l'interligne, au-dessus de « celle des philosophes »
biffé.} Partout la tendance à la généralisation, \struck{produite par} dont
la cause première est\note{« dont la cause première est » ajouté dans
l'interligne.} le sentiment intime de l'unité de l'être, a précipité le
jugement en avant de l'expérience, dont il auroit dû attendre les
\struck{rapports} données.

\note{À la fin de ce paragraphe, un nom d'une grande écriture, souligné,
lu \uncertain{Pemeguin} ; à côté, au crayon et très pâle, une mention
commençant par « Ch. 5 », dont la suite n'a pas été lue.}
Il nous reste à présenter quelques considérations sur l'état des lettres aux
diverses époques dont nous avons examiné les opinions systématiques.

On sait que les anciens, si mal informés des phénomènes naturels, si
ignorants à l'égard des lois qui régissent les faits qu'ils ne pouvoient pas
ignorer, et si fertiles pourtant en généralités propres à embrasser l'univers
entier, avoient atteint la perfection dans tous les genres d'écrire. Ne nous
en étonnons pas. Le sentiment du beau étoit pris dans \struck{la nature} les
lois mêmes\note{« les lois mêmes » ajouté dans l'interligne.} de la nature
intellectuelle de l'homme ; les observations n'avoient besoin ni du secours
des instrumens imaginés par les modernes, ni de la constance et de la
maturité de raison, qui, chez \struck{eux} ces derniers semble avoir remplacé
cette fraîcheur d'imagination qui \struck{devoit} dut peut-être à son entière
indépendance ce qu'elle eut de force et de grâce.

\page{80}
\struck{Lorsque l'étude se bornoit} L'art d'émouvoir et celui de plaire n'ont
besoin que de la connoissance des choses humaines.\note{Un crochet dans la
marge, devant « L'art », comme aux vues 61 et 70.} Il étoit dans la nature
\struck{des choses} de l'esprit humain, de se réfléchir d'abord vers lui
même. S'il a pu errer, en y cherchant le modèle de l'univers, et le but, la
cause finale, de toutes les existences placées en dehors de la sienne, il ne
pouvoit se tromper \struck{à l'égard des} par rapport aux lois de son être.
À cet égard, l'homme s'est trouvé naturellement placé dans la position qu'il
n'a prise que fort tard par rapport aux objets extérieurs. Il a observé les
faits intellectuels ; ils étoient trop près de lui pour qu'il ne sût pas les
bien voir.

Le génie, qui sait à son gré \struck{transmettre} reproduire et transmettre
des impressions profondes, ne pouvoit \struck{attendre} manquer de manifester
sa puissance, aussitôt que l'homme, placé dans l'état social, s'est trouvé
environné de ses semblables.

Sans doute, le goût est le fruit d'un grand nombre d'observations, et il n'a
pu être fixé que \struck{par la} longtems après l'apparition des
\struck{\ill{}} premiers ouvrages qui en offroient le modèle. Mais enfin,
quelque soit la variété des genres, \uncertain{on sait} pourtant qu'il ne
pouvoit pas se passer un tems immense avant que les \struck{réflexions}
observations, les remarques, et les comparaisons se trouvassent assez
multipliées pour avoir fourni à l'intelligence humaine tout ce qu'elle est
susceptible d'acquérir, dans un genre d'étude \struck{qui} exempt\note{« exempt »
dans la marge de gauche.}, par sa nature, \struck{étoit exempt} des causes
d'erreurs qui \uncertain{l'avoient} égaré dans des recherches où l'objet de

\end{document}
